\section{Method}
\label{sec:method}

OPENER is built for the legal domain, yet none of its components is fitted to it. The detector and the entity embedder are ready-to-use checkpoints kept frozen, so no gradient step is ever taken on legal text. Only the typing head touches the legal side: OPENER-ZS consumes no legal label at all, deriving its prototypes from the type names and refining them on the unlabelled test mentions, while OPENER-Sup fits a linear classifier on the frozen embeddings of the target's training mentions, consuming legal labels without updating a single encoder weight. We therefore measure how far a frugal open-world pipeline goes on legal NER \emph{without} being retrained on it, which is the regime a legal deployment actually faces.

\begin{figure*}[t]
\centering
\IfFileExists{assets/opener-architecture-v3.pdf}{%
  \includegraphics[width=0.52\linewidth]{assets/opener-architecture-v3.pdf}%
}{\IfFileExists{../assets/opener-architecture-v3.pdf}{%
  \includegraphics[width=0.52\linewidth]{../assets/opener-architecture-v3.pdf}%
}{\fbox{\parbox{0.9\linewidth}{\centering architecture figure (opener-architecture-v3.pdf)}}}}
\caption{The OPENER pipeline (Detect\,$\rightarrow$\,Embed\,$\rightarrow$\,Type), with no component fitted to legal data, with its two typing heads: OPENER-ZS (label-name prototypes, transductive refinement, detector fusion) and OPENER-Sup (supervised linear probe). Nothing is retrained on legal data. The legal AMI figures shown are means over the three embedder-retraining seeds. The end-to-end ones use the full OPENER-ZS head, the typing-on-gold ones its first step alone, since no detector is in the loop (\autoref{sec:exp:setup}).}
\label{fig:pipeline}
\end{figure*}

\subsection{\rev{Overview and notation}}
\label{sec:method:overview}
\rev{The pipeline runs in three stages, and \autoref{fig:pipeline} shows them in order.} The detector maps a sentence $s$ to a set of candidate mentions $\mathcal{M}(s)=\{m_i\}$, each mention is embedded \emph{in context}, and a typing head assigns it a type drawn from the target inventory $\mathcal{T}$,
\begin{equation}
e_i = f_\theta(s, m_i)\in\mathbb{R}^{D},\qquad \hat y_i\in\mathcal{T},
\end{equation}
where $f_\theta$ is the frozen embedder and $e_i$ the $D$-dimensional vector it produces for mention $m_i$, read together with its sentence $s$. We keep the full $D{=}768$. The predicted type $\hat y_i$ is drawn from $\mathcal{T}$, the legal type inventory of the target corpus, which $f_\theta$ never saw during its own training. \rev{Stages one and two are identical at both operating points. Only the typing head of stage three changes, and \autoref{tab:protocols} in \autoref{sec:appendix:protocols} states exactly how.}

\subsection{\rev{Preprocessing}}
\label{sec:method:preprocessing}
\rev{Three preparation steps sit between a raw legal corpus and the pipeline, and each one is a deliberate choice rather than a formatting detail.}

\rev{\paragraph{Corpus normalisation} The four corpora ship in different formats (CoNLL columns, character spans, token-level tags). We convert them all to one representation, a sentence paired with $(\text{start},\text{end},\text{type})$ triples in \emph{character} offsets, so a single evaluation path serves every corpus and no system benefits from a corpus-specific reader. Sentences with no gold entity are dropped, on every system alike, and each test split is capped at its first $1000$ entity-bearing sentences.}

\rev{\paragraph{Span-in-context encoding} The embedder never sees a mention in isolation. A mention spanning characters $[a,b)$ of sentence $s$ is encoded as the whole sentence with the mention delimited in place,}
\rev{\begin{equation*}
\texttt{classification:}\;s_{<a}\;\texttt{[ENT]}\;s_{a:b}\;\texttt{[/ENT]}\;s_{\geq b},
\end{equation*}}
\rev{where \texttt{classification:} is the task prefix Nomic v1.5 expects. Keeping the sentence is what lets the same surface form be typed differently in two contexts, which legal text demands constantly (a party name that is a company in one clause and a litigant in the next). The markers, rather than a cropped span, keep that context while still telling the encoder which tokens to describe.}

\rev{\paragraph{Type anchors} OPENER-ZS needs a vector per type, and a raw schema code such as \texttt{ORGANIZACAO} or \texttt{GS} carries little meaning for an English-trained encoder. Each type therefore receives one to five English anchor words (\texttt{COURT} $\rightarrow$ \emph{court}, \emph{tribunal}, \emph{judiciary}), listed in full in \autoref{sec:appendix:setup} and fixed before any result was seen. End-to-end, every anchor is embedded through the same four templates, \emph{\{a\}}, \emph{a \{a\}}, \emph{the \{a\}} and \emph{an example of \{a\}}, each one passed through the span-in-context format above with the anchor itself marked. The point is distributional: a prototype built this way is produced by the same input format as the mentions it will be compared against, instead of a bare word the encoder would see in a format it never meets at inference.}

\subsection{\rev{Mention detection}}
\label{sec:method:detection}
Spans are proposed by a frozen GLiNER detector \citep{zaratiana2024gliner} \rev{prompted with the target type names at confidence threshold $0.3$, the value released with the model. The two heads run on different sizes, GLiNER-L for OPENER-ZS and GLiNER-M for OPENER-Sup, because the detector is a modular slot rather than a tuned hyperparameter. \autoref{tab:legal-main} reports every GLiNER size on its own, so the effect of that choice stays readable from the baselines rather than hidden inside our pipeline.}

\subsection{\rev{Entity embedding}}
\label{sec:method:embedding}
Mentions are embedded by Nomic v1.5 \citep{nussbaum2024nomic}, a Sentence-BERT-style \citep{reimers2019sbert}, Matryoshka-truncatable embedder \citep{kusupati2022matryoshka} that we sharpen once, with a FaceNet triplet step \citep{schroff2015facenet} on English, general-domain CoNLL-2003 \citep{tjongkimsang2003conll} spans and an error-driven hard-negative mining pass. Training it away from any legal schema is deliberate, and the same single checkpoint then serves all four legal schemas here without re-fitting. Nothing in the pipeline is legal-specific except the type inventory it is asked about, and that is exactly the design this paper puts to the test. \rev{We keep the full $768$ dimensions and use no Matryoshka truncation, since the embedder is not the cost bottleneck of the pipeline (\autoref{tab:legal-main}) and truncating would trade quality for a saving we do not need.}
\rev{Both embedder checkpoints are public.\footnote{\scriptsize\url{https://huggingface.co/Thibault-GAREL/opener-sup} and \url{https://huggingface.co/Thibault-GAREL/opener-zs}.}}

\subsection{Entity typing}
\label{sec:method:typing}
\paragraph{OPENER-ZS} The zero-shot head, the open-world default, uses no target labels. Each legal type $t$ is represented by a prototype $p_t$, the $\ell_2$-normalised centroid of the embeddings of the words forming its name,
\begin{equation}
\begin{split}
\overline{e}_t&=\frac{1}{|\mathrm{name}(t)|}\!\!\sum_{w\,\in\,\mathrm{name}(t)}\!\!f_\theta(w),\\[2pt]
p_t &= \overline{e}_t\,\big/\,\lVert \overline{e}_t\rVert_2,
\end{split}
\end{equation}
where $\mathrm{name}(t)$ is the set of anchor words naming type $t$ (\autoref{sec:method:preprocessing}). Each anchor word $w$ contributes its embedding $f_\theta(w)$, and $\overline{e}_t$ is their average. Dividing that average by its Euclidean norm $\lVert\overline{e}_t\rVert_2$ gives a unit-length prototype. \rev{End-to-end, $f_\theta(w)$ is replaced by the four-template average of \autoref{sec:method:preprocessing}, which changes the vector but not the equation.}

A mention is assigned the nearest prototype by cosine similarity, refined by a \emph{transductive} pass of spherical $k$-means over the unlabelled test mentions, seeded by the prototypes, and \emph{fused} with the detector's own zero-shot label,
\begin{equation}
s_t(i)=\cos(e_i,\tilde{p}_t)+\beta\,g_i\,\indic\!\left[t=\ell_{\text{GLiNER}}(i)\right],
\end{equation}
where $s_t(i)$ is the fused score of assigning type $t$ to mention $m_i$, and $\tilde{p}_t$ the centroid of type $t$ after that refinement, which is $p_t$ itself when the transductive pass is switched off. The detector contributes the type $\ell_{\text{GLiNER}}(i)$ it predicted for that span and its confidence $g_i\in[0,1]$ in it. The indicator $\indic[\cdot]$ is worth $1$ only when $t$ is that predicted type, so the fusion term only ever raises the score of the detector's own choice, and the mention is typed as $\hat y_i=\arg\max_t s_t(i)$. \rev{Two constants govern this head, and both are fixed a priori rather than tuned on a legal set. The refinement runs three iterations, enough for the centroids to settle on corpora of this size and few enough to keep them anchored to the label names that seeded them. The fusion weight $\beta{=}0.05$ is small on purpose, so the detector can only break ties the prototypes leave open instead of overriding them, a bound \autoref{sec:appendix:beta} then sweeps post-hoc.}

\paragraph{OPENER-Sup} A balanced linear classifier (scikit-learn's LinearSVC \citep{pedregosa2011scikit} with \texttt{class\_weight=balanced}, built on LIBLINEAR \citep{fan2008liblinear}) is fit one-vs-rest on the legal target's training labels. For each type $c$, it minimises a squared-hinge objective with balanced class weights,
\begin{equation}
\begin{split}
\min_{w_c,b_c}\;&\tfrac12\lVert w_c\rVert_2^2\\
&+\,C\sum_i \omega_{y_i}\,\max\!\big(0,\,1-z_{ic}\,h_c(e_i)\big)^2,
\end{split}
\end{equation}
with $h_c(e_i)=w_c^{\top} e_i+b_c$ the one-vs-rest hyperplane of type $c$, $e_i$ the frozen embedding of training mention $i$, $y_i$ its gold legal type, $z_{ic}={+}1$ if $y_i{=}c$ and ${-}1$ otherwise, and $C$ the regularisation strength. The balancing weight $\omega_c=n/(|\mathcal{T}|\,n_c)$ up-weights rare types, with $n$ the total number of training mentions and $n_c$ the count of type $c$. This matters on legal schemas, where a single type can hold well under $1\%$ of the mentions (Court in E-NER, \autoref{tab:legal-datasets}). A mention is typed as $\hat y_i=\arg\max_c h_c(e_i)$. Rather than a single label-name match, this carves the embedding space into one region per type, separating types whose vectors sit close together as long as their regions do not overlap. This is the supervised upper bound, and it needs target labels. \rev{End-to-end the probe is fit on the detector's own training spans rather than on gold ones, so that it learns to type the noisy boundaries it will actually be handed at inference.}

\subsection{\rev{Evaluation protocol}}
\label{sec:method:evaluation}
We report both \emph{typing-on-gold} (perfect detection) and \emph{end-to-end} protocols. Predicted and gold mentions are aligned by exact character offsets, and detection errors stay in the evaluation through sentinel labels, one for a gold mention no system predicted and one for a predicted span with no gold counterpart, following the alignment protocol of \citet{genest2025owner}. Sentences carrying no gold entity are not evaluated, on any system.

\paragraph{\rev{Quality}} Quality is the Adjusted Mutual Information (AMI) \citep{vinh2010ami} between the gold type partition $U$ and the predicted partition $V$ over the aligned mentions,
\begin{equation}
\mathrm{AMI}=\frac{\mathrm{MI}-\mathbb{E}\!\left[\mathrm{MI}\right]}{\tfrac12\!\left[H(U)+H(V)\right]-\mathbb{E}\!\left[\mathrm{MI}\right]},
\end{equation}
where $\mathrm{MI}=\mathrm{MI}(U,V)$ is the mutual information between the two partitions, which measures how much a mention's predicted type reduces the uncertainty about its gold type, and $H$ the entropy. Subtracting the chance agreement $\mathbb{E}[\mathrm{MI}]$ makes a random typing score near $0$ whatever the number of types, a perfect one $1$. The denominator uses the arithmetic-mean normalisation (the scikit-learn default). \rev{We report it as AMI\,$\times100$, so the scale runs from ${\approx}0$ for chance to $100$ for a partition identical to the gold one. The value $100$ is reachable in the typing-on-gold protocol, where both partitions cover the same mentions, but end-to-end it is not: a missed or spurious span puts a sentinel on one side and a real type on the other, so any detection error caps the score below $100$ whatever the typing head does. \autoref{sec:exp:results} quantifies that cap directly.}

\rev{\paragraph{Why AMI, and what we report beside it} AMI is invariant to a relabelling of the predicted types: a system that consistently called \emph{judges} what the gold calls \emph{lawyers}, and conversely, would score $100$ AMI and $0$ F1. That invariance is the point here, since OWNER names its clusters dynamically and the generative baselines paraphrase type strings, so a metric requiring predicted and gold names to coincide cannot score them at all. AMI is also corrected for chance, which matters across schemas of $6$ to $19$ unequally frequent types. The argument binds only for systems answering outside the target inventory, though. OPENER's heads, the GLiNER family, GNER and the LLM baseline all answer inside it, so for these an F1 is well defined and we report it: macro-F1 on gold mentions (\autoref{tab:legal-gold}) and micro entity-level precision, recall and F1 end-to-end (\autoref{tab:legal-decomp}), a span counting as correct only when offsets and type both match, the convention of the four dataset papers. OWNER alone has none.}

\paragraph{\rev{Cost}} We also report per-sentence latency (p50) and energy (Wh per $1000$ sentences, measured with CodeCarbon \citep{codecarbon2024}).
